
%\needspace{4\baselineskip}
\color{uniblue}{\subsection{Аварийная максимальная токовая защита (АМТЗ)}}
\color{black}

\begin{enumerate}[label=\arabic{section}.\arabic{subsection}.\arabic{enumi}, labelsep=4pt, leftmargin=0pt, itemindent=57pt, itemsep=0pt, parsep=0pt] % если что можно крутить itemsep
%%%%%%%%%%%%%%%%%%%%%%%%%%%%%%%%%%%%%%%%%%%%%%%%%%%%%%%%%%%%%%%%%%%%%%  ОБЩИЕ СВЕДЕНИЯ %%%%%%%%%%%%%%%%%%%%%%%%%%%%%%%%%%%%%%%%%%%%%%%%%%

\item
Аварийная максимальная токовая защита (АМТЗ) применяется в качестве аварийной защиты при неисправностях в цепях напряжения.

\item В составе АМТЗ предусмотрены:
\begin{itemize}
    \item  общая логика;
    \item  2 ступени, реагирующие на величину фазных токов (АМТЗ-Ф);
    \item 2 ступени, реагирующие на величину тока обратной последовательности (АМТЗ-ОП);
    \item 2 ступени, реагирующие на величину тока нулевой последовательности (АМТЗ-НП); 
\end{itemize}

\item
Функциональная схема общей логики АМТЗ приведена на рисунке \ref{pict_amtz:common_logic}.

\vspace{2mm}
\begin{figure}[H]
\centering
\includegraphics[width=0.8\textwidth,height=0.8\textheight,keepaspectratio]{\fbpath/561.0161 АМТЗ/img/common_logic.pdf}
\captionof{figure}{Функциональная схема общей логики АМТЗ}\label{pict_amtz:common_logic}
\end{figure}

\item
Функциональная схема логики АМТЗ-Ф-1(2) приведена на рисунке \ref{pict_amtz:AMTZ}.

\vspace{2mm}
\begin{figure}[H]
\centering
\includegraphics[width=1\textwidth,height=1\textheight,keepaspectratio]{\fbpath/561.0161 АМТЗ/img/AMTZ.pdf}
\captionof{figure}{Функциональная схема ступеней АМТЗ-Ф-1(2)}\label{pict_amtz:AMTZ}
\end{figure}

\item
Уставки АМТЗ представлены в таблице \ref{app_set:amtz}.


\item
АМТЗ-Ф вводятся в работу программными ключами \textbf{«АМТЗ-Ф-1»} и \textbf{«АМТЗ-Ф-2»}. Токи срабатывания задаются уставками \textbf{«Iср АМТЗ-Ф-1»} и \textbf{«Iср АМТЗ-Ф-2»}. Действие ступеней АМТЗ-Ф осуществляется с регулируемыми независимыми выдержками времени \textbf{«Тср АМТЗ-Ф-1»} и \textbf{«Тср АМТЗ-Ф-2»}.

\item
При необходимости можно ввести контроль от БК-I-м отдельно для каждой ступени программными ключами \textbf{«БК-I-м АМТЗ-Ф-1»} и \textbf{«БК-I-м АМТЗ-Ф-2»}.

\item
Ввод в работу ступеней АМТЗ-ОП и АМТЗ-НП, настройка параметров срабатывания, ввод контроля от БК-I-м осуществляется аналогично.

\item
Автоматический ввод в работу АМТЗ при возникновении неисправности в цепях напряжения производится при отсутствии внешнего сигнала \textit{«Вывод АМТЗ»}. Пуски ступеней возможны даже при отсутствии неисправности в цепях напряжения, но выходная логика срабатывания находится в заблокированном состоянии. При появлении входного сигнала \textit{«БНН»} выходная логика ступеней на срабатывание деблокируется, при этом формируется логический сигнал \textit{«АМТЗ активна»}. Предусмотрена также возможность оперативного ввода АМТЗ (независимо от сигнала срабатывания БНН) от внешнего сигнала \textit{«Опер. ввод АМТЗ»}.


\end{enumerate}